\subsubsection*{1.4.3.}
对 $\mathcal D$ 上的拓扑空间 $\mathcal X$，如下赋予 $D$ 上的拓扑空间 $X=\beta^*\mathcal X$。在 $X^*=\beta^*\mathcal X^*$ 上添加 $0$ 上的纤维 $X_0$。$x_0\in X_0$ 的一个基本邻域系由集合
\[
V\cup
\bigl(\mathrm{sp}^{-1}(V)\times\widetilde U^*/\pi_1(D^*)\bigr),
\]
构成，其中 $V$ 为 $X_0$ 中 $x_0$ 的邻域，$U$ 为 $D$ 中 $0$ 的邻域，并且 $\widetilde U^*=U\times_D\widetilde D^*$。
用 $j$（分别用 $i$）表示 $X^*$（分别为 $X_0$）到 $X$ 的包含。对 $\mathcal X^*$ 上的层 $F$，有
\[
i^*j_*\beta^*F
=(\mathrm{sp}_*F_\eta)^{\pi_1(D^*)}.
\]
由此有拓扑斯的交换图
\[
\begin{tikzcd}
X \arrow[r, "\beta"] \arrow[d]
  & \mathcal X \arrow[d] \\
D \arrow[r, "\beta"]
  & \mathcal D
\end{tikzcd}
\]
使对 $\mathcal X$ 上的每个层 $F$，存在同构
\begin{equation}\tag{1.4.3.1}
\Psi(F)\xrightarrow{\ \sim\ }\Psi(\beta^*F).
\end{equation}

\subsubsection*{命题 1.4.4.}
映射 (1.4.3.1) 导出一个从 $D^+(\mathcal X)$ 到 $D^+(X_0\times\mathcal D)$ 的函子之间的同构
\[
R\Psi\xrightarrow{\ \sim\ }R\Psi\circ\beta^*
\]

事实上，设 $\overline j$ 为
\[
\overline X^*=X\times_D\widetilde D^*
\simeq X_\eta\times\widetilde D^*
\]
到 $\overline X=X\times_D\widetilde D$ 的包含。此命题归结为验证：对 $X_\eta$ 上的每个松交换群层 $F$，层 $\operatorname{pr}_1^*F$ 关于 $\overline j_*$ 是无上同调的。这又归结为（取 $B=\widetilde D^*$）下述可用同伦论证证明的熟知断言。若 $F$ 是空间 $Y$ 上的交换群层，$B$ 是圆盘，则
\[
H^i(Y,F)\xrightarrow{\ \sim\ }
H^i(Y\times B,\operatorname{pr}_1^*F).
\]

\subsubsection*{1.4.5.}
命题 1.4.4 给出了对 $X$ 上形如 $\beta^*F$ 的层计算 $R\Psi$ 的一种方法，将 $R\Psi$ 的计算实质上归结为 $R\,\mathrm{sp}_*$ 的计算。

